\chapter{QUBO et modèle d'Ising}\label{ch:qubo}
\iterbanner{V0}

C'est le chapitre charnière : il traduit un problème d'optimisation d'ingénieur en un hamiltonien que le recuit ou QAOA savent traiter.

\section{Optimisation combinatoire et QUBO}
\begin{definition}[ (QUBO)]
Un problème \textbf{QUBO} (\emph{Quadratic Unconstrained Binary Optimization}) consiste à
\[\min_{\mathbf b\in\{0,1\}^n}\ E(\mathbf b)=\mathbf b^\top Q\,\mathbf b+c=\sum_iQ_{ii}b_i+\sum_{i<j}Q_{ij}b_ib_j+c,\]
avec $Q$ triangulaire supérieure. Les termes linéaires sont sur la diagonale, car $b_i^2=b_i$.
\end{definition}
\begin{itemize}
\item \emph{Quadratic} : pas de produit de trois bits ou plus.
\item \emph{Unconstrained} : toutes les $2^n$ chaînes sont admissibles. Les contraintes doivent être transformées en pénalités (§\ref{sec:penalites}).
\item \emph{Binary} : variables 0/1. Les entiers doivent être encodés (§\ref{sec:encodage}).
\end{itemize}
Le QUBO général est \textbf{NP-difficile} : il contient Max-Cut, le voyageur de commerce, l'ordonnancement… On ne connaît pas d'algorithme, classique ou quantique, qui le résolve en temps polynomial dans le pire cas. Recuit et QAOA sont des heuristiques.

\section{Deux exemples}
\paragraph{Le QUBO jouet du cours.} Trois bits sur une chaîne $0-1-2$ ; on veut en activer le plus possible sans que deux voisins soient actifs (« ensemble indépendant maximal ») :
\[E=-b_0-b_1-b_2+2b_0b_1+2b_1b_2,\qquad Q=\begin{pmatrix}-1&2&0\\0&-1&2\\0&0&-1\end{pmatrix}.\]
Chaque bit actif rapporte $-1$, chaque paire de voisins actifs coûte $+2$. Énumération : $E(101)=-2$ est le minimum unique (figure~\ref{fig:niveaux}). Ce problème à 3 qubits sert de fil rouge pour le recuit et QAOA, car tout s'y calcule à la main ou presque.

\paragraph{Max-Cut.} Couper un graphe en deux groupes en maximisant le nombre d'arêtes coupées. Avec $b_i$ le groupe du nœud $i$, l'arête $(i,j)$ est coupée si $b_i\neq b_j$, ce qui s'écrit $b_i+b_j-2b_ib_j=1$. On minimise donc $E=-\sum_{(i,j)}(b_i+b_j-2b_ib_j)$. C'est l'exemple historique de QAOA (Farhi et al., 2014).

\section{Du QUBO au modèle d'Ising}
Le matériel quantique manipule des \emph{spins} $s_i=\pm1$, valeurs propres de $Z_i$, plutôt que des bits. On substitue $b_i=(1-s_i)/2$ :
\begin{definition}[ (modèle d'Ising)]
\[E(\mathbf s)=\sum_ih_is_i+\sum_{i<j}J_{ij}s_is_j+\text{cte},\]
\[J_{ij}=\frac{Q_{ij}}4,\qquad h_i=-\frac{Q_{ii}}2-\frac14\sum_{j\neq i}Q_{\min(i,j)\max(i,j)},\qquad \text{cte}=c+\frac12\sum_iQ_{ii}+\frac14\sum_{i<j}Q_{ij}.\]
L'\textbf{hamiltonien de coût} s'obtient en remplaçant $s_i$ par l'opérateur $Z_i$ :
\[H_C=\sum_ih_iZ_i+\sum_{i<j}J_{ij}Z_iZ_j+\text{cte}\cdot\Id .\]
\end{definition}
\begin{itemize}
\item $h_i$ est un \emph{champ local}, qui favorise un spin particulier.
\item $J_{ij}>0$ favorise des spins opposés (antiferromagnétique) ; $J_{ij}<0$ favorise des spins égaux (ferromagnétique).
\item Le vocabulaire vient du magnétisme. Ising (1925) modélisait des aimants ; c'est le même objet mathématique.
\end{itemize}
Pour le jouet : $h=(0,\,-\tfrac12,\,0)$, $J_{01}=J_{12}=\tfrac12$, constante $-\tfrac12$ (exercice~\ref{ch:qubo}.1).

\begin{figure}[htbp]
\centering
\begin{tikzpicture}[node distance=22mm,sp/.style={circle,draw=cink,thick,minimum size=9mm,font=\small}]
\node[sp] (0) {$s_0$}; \node[sp,right=of 0] (1) {$s_1$}; \node[sp,right=of 1] (2) {$s_2$};
\draw[corange,very thick] (0) -- node[above,font=\small] {$J_{01}=+\frac12$} (1);
\draw[corange,very thick] (1) -- node[above,font=\small] {$J_{12}=+\frac12$} (2);
\node[below=2mm of 0,font=\small] {$h_0=0$}; \node[below=2mm of 1,font=\small] {$h_1=-\frac12$}; \node[below=2mm of 2,font=\small] {$h_2=0$};
\end{tikzpicture}
\caption{Le QUBO jouet comme graphe d'Ising : un nœud par qubit, une arête par couplage non nul. Le nombre d'arêtes fixe le nombre de portes $R_{ZZ}$ (QAOA) ou de coupleurs (recuit quantique).}\label{fig:graphe-jouet}
\end{figure}

\section{Contraintes : les pénalités}\label{sec:penalites}
Un QUBO n'a pas de contraintes. On les ajoute à l'énergie sous forme de \textbf{pénalité}, nulle si la contrainte est satisfaite et positive sinon :
\begin{itemize}
\item \textbf{égalité} $g(\mathbf b)=0$ : ajouter $P\,g(\mathbf b)^2$. Exemple, « exactement un bit actif parmi $k$ » (one-hot) : $P\big(\sum_kb_k-1\big)^2$ ;
\item \textbf{inégalité} $g(\mathbf b)\le0$ avec $g$ entier : introduire des bits d'écart $a_l$ et imposer $g(\mathbf b)+\sum_l2^la_l=0$ par pénalité. Cela coûte des qubits supplémentaires.
\end{itemize}
\begin{piege}[ (degré)]
La pénalité est un \emph{carré}. Si $g$ est linéaire en bits, $g^2$ est quadratique : on reste dans le QUBO. Si $g$ est déjà quadratique (par exemple $C_L$ approché par un modèle quadratique), $g^2$ est de \textbf{degré 4} : ce n'est plus un QUBO. Il faut alors linéariser $g$, ou reformuler le problème pour que la contrainte disparaisse.
\end{piege}
\paragraph{Réglage de $P$.} $P$ doit être assez grand pour qu'aucune violation ne soit rentable, mais pas trop : un $P$ énorme écrase les écarts d'énergie utiles. QAOA devient alors moins sensible au coût, et le recuit quantique perd en précision (coefficients bornés et quantifiés sur le matériel).

\section{Encoder des entiers en bits}\label{sec:encodage}
Une variable entière $x\in\{0,\dots,L-1\}$ peut s'écrire de plusieurs façons :
\begin{center}
\begin{tabular}{lll c}
\toprule
Encodage & $x=$ & contrainte & bits pour $L=8$\\
\midrule
binaire & $\sum_{j}2^jb_j$ & aucune ($L=2^k$) & 3\\
one-hot & $\sum_jj\,b_j$ & $\sum_jb_j=1$ (pénalité) & 8\\
unaire (« thermomètre ») & $\sum_jb_j$ & ordre $b_{j+1}\le b_j$, ou aucune & 7\\
\bottomrule
\end{tabular}
\end{center}
Le binaire est le plus économe en qubits, mais deux niveaux voisins peuvent différer de plusieurs bits : passer de 3 (011) à 4 (100) en change trois. Cela rend le paysage d'énergie plus « rugueux » pour un algorithme qui modifie un bit à la fois, comme le recuit simulé.

\section{Notre QUBO : du modèle de traînée à la matrice $Q$}\label{sec:notre-qubo}
\paragraph{Variables.} $m$, $p$, $t$ prennent chacune 8 niveaux d'indice $x_i\in\{0,\dots,7\}$, codés en binaire sur 3 bits : $x_i=b_{i0}+2b_{i1}+4b_{i2}$. Cela fait 9 bits au total, donc \textbf{9 qubits}.

\paragraph{Modèle.} On ajuste par moindres carrés, sur des profils simulés par la phase~1, un polynôme quadratique en les variables normalisées $z_i=2x_i/7-1\in[-1,1]$ :
\[\widehat{C_D}(z)=c_0+\mathbf g^\top\mathbf z+\mathbf z^\top M\mathbf z\qquad(1+3+6=10\text{ coefficients}).\]

\paragraph{Passage en bits.} Chaque $z_i$ est \emph{affine} en bits, $\mathbf z=\mathbf a_0+A\mathbf b$, avec $\mathbf a_0=-\mathbf 1$ et $A_{i,(i,j)}=2\cdot2^j/7$. En substituant :
\[\widehat{C_D}=\underbrace{c_0+\mathbf g^\top\mathbf a_0+\mathbf a_0^\top M\mathbf a_0}_{c}+\underbrace{(\mathbf g+2M\mathbf a_0)^\top A\,\mathbf b}_{\text{linéaire}\to\text{diag}(Q)}+\mathbf b^\top\underbrace{(A^\top MA)}_{\text{quadratique}}\mathbf b .\]
Un polynôme de degré 2 en $z$ reste de degré 2 en bits : \textbf{c'est un QUBO, sans approximation}. Le code \texttt{quantique/qubo.py} (\texttt{surrogate\_to\_qubo}) le vérifie : énergie QUBO et prédiction du modèle coïncident à $10^{-14}$ près sur les 512 profils.

\begin{figure}[htbp]
\centering
% figure générée par tools/make_figures.py : matrice_Q
\begin{tikzpicture}[scale=0.62]
\fill[corange!10] (0,0) rectangle ++(1,-1);
\node[font=\tiny] at (0.5,-0.5) {1.3};
\fill[corange!14] (1,0) rectangle ++(1,-1);
\node[font=\tiny] at (1.5,-0.5) {1.8};
\fill[corange!28] (2,0) rectangle ++(1,-1);
\node[font=\tiny] at (2.5,-0.5) {3.6};
\fill[cblue!4] (3,0) rectangle ++(1,-1);
\node[font=\tiny] at (3.5,-0.5) {-0.5};
\fill[cblue!8] (4,0) rectangle ++(1,-1);
\node[font=\tiny] at (4.5,-0.5) {-1.1};
\fill[cblue!17] (5,0) rectangle ++(1,-1);
\node[font=\tiny] at (5.5,-0.5) {-2.1};
\fill[corange!2] (6,0) rectangle ++(1,-1);
\node[font=\tiny] at (6.5,-0.5) {0.2};
\fill[corange!4] (7,0) rectangle ++(1,-1);
\node[font=\tiny] at (7.5,-0.5) {0.5};
\fill[corange!8] (8,0) rectangle ++(1,-1);
\node[font=\tiny] at (8.5,-0.5) {1.0};
\fill[corange!28] (1,-1) rectangle ++(1,-1);
\node[font=\tiny] at (1.5,-1.5) {3.6};
\fill[corange!55] (2,-1) rectangle ++(1,-1);
\node[font=\tiny] at (2.5,-1.5) {7.1};
\fill[cblue!8] (3,-1) rectangle ++(1,-1);
\node[font=\tiny] at (3.5,-1.5) {-1.1};
\fill[cblue!17] (4,-1) rectangle ++(1,-1);
\node[font=\tiny] at (4.5,-1.5) {-2.1};
\fill[cblue!33] (5,-1) rectangle ++(1,-1);
\node[font=\tiny] at (5.5,-1.5) {-4.3};
\fill[corange!4] (6,-1) rectangle ++(1,-1);
\node[font=\tiny] at (6.5,-1.5) {0.5};
\fill[corange!8] (7,-1) rectangle ++(1,-1);
\node[font=\tiny] at (7.5,-1.5) {1.0};
\fill[corange!15] (8,-1) rectangle ++(1,-1);
\node[font=\tiny] at (8.5,-1.5) {2.0};
\fill[corange!83] (2,-2) rectangle ++(1,-1);
\node[font=\tiny] at (2.5,-2.5) {10.7};
\fill[cblue!17] (3,-2) rectangle ++(1,-1);
\node[font=\tiny] at (3.5,-2.5) {-2.1};
\fill[cblue!33] (4,-2) rectangle ++(1,-1);
\node[font=\tiny] at (4.5,-2.5) {-4.3};
\fill[cblue!66] (5,-2) rectangle ++(1,-1);
\node[font=\tiny] at (5.5,-2.5) {-8.5};
\fill[corange!8] (6,-2) rectangle ++(1,-1);
\node[font=\tiny] at (6.5,-2.5) {1.0};
\fill[corange!15] (7,-2) rectangle ++(1,-1);
\node[font=\tiny] at (7.5,-2.5) {2.0};
\fill[corange!31] (8,-2) rectangle ++(1,-1);
\node[font=\tiny] at (8.5,-2.5) {3.9};
\fill[cblue!7] (3,-3) rectangle ++(1,-1);
\node[font=\tiny] at (3.5,-3.5) {-0.9};
\fill[corange!15] (4,-3) rectangle ++(1,-1);
\node[font=\tiny] at (4.5,-3.5) {1.9};
\fill[corange!29] (5,-3) rectangle ++(1,-1);
\node[font=\tiny] at (5.5,-3.5) {3.8};
\fill[cblue!1] (6,-3) rectangle ++(1,-1);
\node[font=\tiny] at (6.5,-3.5) {-0.2};
\fill[cblue!3] (7,-3) rectangle ++(1,-1);
\node[font=\tiny] at (7.5,-3.5) {-0.3};
\fill[cblue!5] (8,-3) rectangle ++(1,-1);
\node[font=\tiny] at (8.5,-3.5) {-0.7};
\fill[cblue!6] (4,-4) rectangle ++(1,-1);
\node[font=\tiny] at (4.5,-4.5) {-0.8};
\fill[corange!59] (5,-4) rectangle ++(1,-1);
\node[font=\tiny] at (5.5,-4.5) {7.6};
\fill[cblue!3] (6,-4) rectangle ++(1,-1);
\node[font=\tiny] at (6.5,-4.5) {-0.3};
\fill[cblue!5] (7,-4) rectangle ++(1,-1);
\node[font=\tiny] at (7.5,-4.5) {-0.7};
\fill[cblue!11] (8,-4) rectangle ++(1,-1);
\node[font=\tiny] at (8.5,-4.5) {-1.4};
\fill[corange!17] (5,-5) rectangle ++(1,-1);
\node[font=\tiny] at (5.5,-5.5) {2.1};
\fill[cblue!5] (6,-5) rectangle ++(1,-1);
\node[font=\tiny] at (6.5,-5.5) {-0.7};
\fill[cblue!11] (7,-5) rectangle ++(1,-1);
\node[font=\tiny] at (7.5,-5.5) {-1.4};
\fill[cblue!22] (8,-5) rectangle ++(1,-1);
\node[font=\tiny] at (8.5,-5.5) {-2.8};
\fill[corange!22] (6,-6) rectangle ++(1,-1);
\node[font=\tiny] at (6.5,-6.5) {2.8};
\fill[cblue!1] (7,-6) rectangle ++(1,-1);
\node[font=\tiny] at (7.5,-6.5) {-0.1};
\fill[cblue!1] (8,-6) rectangle ++(1,-1);
\node[font=\tiny] at (8.5,-6.5) {-0.2};
\fill[corange!43] (7,-7) rectangle ++(1,-1);
\node[font=\tiny] at (7.5,-7.5) {5.5};
\fill[cblue!3] (8,-7) rectangle ++(1,-1);
\node[font=\tiny] at (8.5,-7.5) {-0.3};
\fill[corange!85] (8,-8) rectangle ++(1,-1);
\node[font=\tiny] at (8.5,-8.5) {10.9};
\draw[cmuted,thin] (0,0) grid (9,-9);
\draw[cink,thick] (0,0) rectangle (3,-3) (3,-3) rectangle (6,-6) (6,-6) rectangle (9,-9);
\node[font=\scriptsize] at (0.5,0.35) {$m_0$};
\node[font=\scriptsize] at (1.5,0.35) {$m_1$};
\node[font=\scriptsize] at (2.5,0.35) {$m_2$};
\node[font=\scriptsize] at (3.5,0.35) {$p_0$};
\node[font=\scriptsize] at (4.5,0.35) {$p_1$};
\node[font=\scriptsize] at (5.5,0.35) {$p_2$};
\node[font=\scriptsize] at (6.5,0.35) {$t_0$};
\node[font=\scriptsize] at (7.5,0.35) {$t_1$};
\node[font=\scriptsize] at (8.5,0.35) {$t_2$};
\node[font=\scriptsize] at (-0.45,-0.5) {$m_0$};
\node[font=\scriptsize] at (-0.45,-1.5) {$m_1$};
\node[font=\scriptsize] at (-0.45,-2.5) {$m_2$};
\node[font=\scriptsize] at (-0.45,-3.5) {$p_0$};
\node[font=\scriptsize] at (-0.45,-4.5) {$p_1$};
\node[font=\scriptsize] at (-0.45,-5.5) {$p_2$};
\node[font=\scriptsize] at (-0.45,-6.5) {$t_0$};
\node[font=\scriptsize] at (-0.45,-7.5) {$t_1$};
\node[font=\scriptsize] at (-0.45,-8.5) {$t_2$};
\end{tikzpicture}

\caption{Matrice $Q$ (en counts de traînée) du QUBO aéro, modèle ajusté sur les 512 profils à $\alpha=2^\circ$, $Re=3\cdot10^6$. Bleu : négatif, orange : positif. Les blocs encadrés couplent les 3 bits d'une même variable (termes $z_i^2$) ; les blocs hors diagonale viennent des termes croisés $z_iz_j$. Tous les coefficients hors diagonale sont non nuls : le graphe d'Ising est \textbf{complet} (36 arêtes).}\label{fig:matrice-q}
\end{figure}

\begin{projet}[ : ce que dit le QUBO, et ce qu'il ne dit pas]
\begin{itemize}
\item Minimum du QUBO : $57{,}16$ counts pour le profil $(m,p,t)=(0,\ 0{,}25,\ 6\,\%)$. Le second niveau, $57{,}22$, est le profil $(0,\ 0{,}30,\ 6\,\%)$. Ces deux profils sont \emph{physiquement identiques}, puisque $p$ n'a pas d'effet quand $m=0$ ! Le modèle les sépare de $0{,}06$ count seulement : le fondamental est \emph{quasi dégénéré}.
\item Le vrai $C_D$ de ce profil vaut $61{,}1$ counts : il n'est que \textbf{38\textsuperscript{e} sur 512}. Le vrai minimum, $50{,}5$, est ailleurs. Le QUBO résout \emph{exactement} un problème \emph{approché} : l'erreur vient du modèle quadratique, pas de l'optimiseur. C'est pourquoi la note d'intention prévoit une \textbf{boucle itérative} qui simule les candidats et réajuste le modèle (ch.~\ref{ch:projet}).
\item Le graphe complet a une conséquence matérielle : 36 portes $R_{ZZ}$ par couche de QAOA, et un plongement coûteux sur un recuit quantique (ch.~\ref{ch:recuit}).
\end{itemize}
\end{projet}

\exercices
\begin{exercice}
Convertir le QUBO jouet en modèle d'Ising : calculer $h$, $J$ et la constante avec les formules du cours. Vérifier que $E(\mathbf s)=-2$ pour la chaîne $101$, c'est-à-dire $\mathbf s=(-1,+1,-1)$.
\end{exercice}
\begin{exercice}
On encode $m$, $p$, $t$ (8 niveaux chacun) en one-hot. Combien faut-il de qubits ? Écrire la pénalité qui impose « un seul niveau actif » pour $m$, et dire combien de couplages $Z_iZ_j$ elle ajoute. Conclure sur le choix d'encodage pour notre projet.
\end{exercice}
